% Two-unit repairable redundant system: the three-state birth-death
% graph and the unavailability curve as a function of the
% failure-to-repair ratio r, for one and two repair crews.
\begin{figure}[htbp]
\centering
\begin{tikzpicture}[
  state/.style={draw=engblack, circle, minimum size=1.05cm,
    align=center, font=\scriptsize\bfseries, fill=engblue!12},
  down/.style={draw=engblack, circle, minimum size=1.05cm,
    align=center, font=\scriptsize\bfseries, fill=engred!18},
  edge/.style={-{Stealth}, thick, draw=engblack!80, font=\scriptsize},
]

% ---- Left: birth-death state graph ----
\begin{scope}[shift={(0, 0)}]
  \node[state] (S0) at (0, 0) {$0$\\up-up};
  \node[state] (S1) at (3, 0) {$1$\\one up};
  \node[down]  (S2) at (6, 0) {$2$\\down};

  \draw[edge, bend left=22] (S0) to
    node[above, font=\scriptsize] {$2\lambda$} (S1);
  \draw[edge, bend left=22] (S1) to
    node[below, font=\scriptsize] {$\mu$} (S0);
  \draw[edge, bend left=22] (S1) to
    node[above, font=\scriptsize] {$\lambda$} (S2);
  \draw[edge, bend left=22] (S2) to
    node[below, font=\scriptsize] {$\mu$} (S1);

  \node[anchor=north, font=\small\bfseries] at (3, -1.4)
    {Birth-death graph (single crew)};
  \node[anchor=north, font=\scriptsize] at (3, -2.0)
    {failures = births, repairs = deaths};
\end{scope}

% ---- Right: unavailability vs failure-to-repair ratio ----
\begin{scope}[shift={(9.0, -2.0)}]
  % axes
  \draw[->, thick] (0, 0) -- (4.6, 0)
    node[right, font=\scriptsize] {ratio $r=\lambda/\mu$};
  \draw[->, thick] (0, 0) -- (0, 3.6)
    node[above, font=\scriptsize] {$\bar A$ (unavailability)};

  % single-crew curve  Abar1 = 2r^2/(1+2r+2r^2), scaled y = Abar*7.5
  % points (r, Abar1): (0.1,0.0164) (0.3,0.097) (0.5,0.20) (0.7,0.286)
  %   (0.9,0.355) (1.0,0.40); x = r*4, y = Abar1*7.5
  \draw[thick, engred]
    (0.0, 0.0) -- (0.40, 0.123) -- (1.20, 0.728) -- (2.00, 1.50)
    -- (2.80, 2.14) -- (3.60, 2.66) -- (4.00, 3.00);
  \node[anchor=west, font=\scriptsize, engred] at (4.0, 3.0)
    {single crew};

  % two-crew curve Abar2 = r^2/(1+r)^2, scaled the same way
  % points (r, Abar2): (0.1,0.0083) (0.3,0.053) (0.5,0.111) (0.7,0.169)
  %   (0.9,0.224) (1.0,0.25); x=r*4, y=Abar2*7.5
  \draw[thick, engblue, dashed]
    (0.0, 0.0) -- (0.40, 0.062) -- (1.20, 0.398) -- (2.00, 0.833)
    -- (2.80, 1.27) -- (3.60, 1.68) -- (4.00, 1.875);
  \node[anchor=west, font=\scriptsize, engblue] at (4.0, 1.875)
    {two crews};

  \node[anchor=north, font=\small\bfseries] at (2.3, -0.5)
    {Unavailability vs $r$};
  \node[anchor=north, font=\scriptsize, align=center] at (2.3, -1.0)
    {$\bar A \sim r^{2}$ for small $r$:\\[-1pt]
     redundancy squares the downtime};
\end{scope}

\end{tikzpicture}
\caption[Availability of a two-unit repairable system]{The repairable
redundant system of section~\ref{sec:vol02:ch13:availability}. Left: the
three-state birth-death graph on the number of failed units, with failure
rates $2\lambda$ and $\lambda$ as births and the repair rate $\mu$ as the
death. State $2$ (shaded) is the system-down state. Right: the
steady-state unavailability $\bar A$ as a function of the
failure-to-repair ratio $r = \lambda/\mu$, for one repair crew
($\bar A = 2r^{2}/(1 + 2r + 2r^{2})$, solid) and two crews
($\bar A = r^{2}/(1 + r)^{2}$, dashed). Both curves grow as $r^{2}$ for
small $r$, the quadratic that is the quantitative case for active
redundancy with repair; the second crew halves the unavailability but
matters less than fast first-failure response because state $2$ is
already rare. Source: birth-death solution,
\cite[ch.~6]{text:wasserman2004}.}
\label{fig:vol02:ch13:availability}
\end{figure}
